\chapter{Description of Algorithms}\label{cha:descr-algor}
This chapter summarizes the main features of QDOAS and describes the
structure of the program and the algorithms.

\section{Differential Optical Absorption Spectroscopy}\label{sec:diff-optic-absorpt}
\gls{doas} is a widely used inversion method for the retrieval of
atmospheric trace gas abundances from multi-wavelength light
measurements.  It uses the structured absorption of many trace gases
in the UV, visible and near-infrared spectral ranges. The DOAS method
was originally developed for ground-based measurements
(\citet{platt1994,platt_2008}) and has been successfully adapted to
nadir measurements from UV-Vis space-borne spectrometers
(\citet{sciamachy}). It relies on the application of the Beer-Lambert
law to the whole atmosphere in a limited range of wavelengths.

The Beer-Lambert law states that the radiant intensity traversing a
homogeneous medium decreases exponentially with the product of the
extinction coefficient and the path length.  Applying this law to the
atmosphere, we obtain
\begin{equation}\label{eq:beer-lambert}
  I(\lambda)= I_0(\lambda)\exp \left( -\sum_{j=1}^{n} S^j(\lambda) c_j\right),
\end{equation}
where
\begin{itemize}
\item[$I_0$] is the spectrum at the top of the atmosphere, without extinction;
\item[$I$] is the measured spectrum after extinction in the atmosphere;
\item[$S^j$] is the absorption cross section of the species $j$, with
  wavelength dependent structures [cm2/molec.];
\item[$c_j$] is the column density of the species $j$ [molec./cm2].
\end{itemize}
The logarithm of the ratio of the spectrum $I_0$ (also called the
control spectrum) and the measured spectrum $I$ is denoted optical
density (or optical thickness) $\tau$
\begin{equation}\label{eq:optical_density}
\tau  \equiv \ln\left( \frac{I_0(\lambda)}{I(\lambda)}\right)= \sum_{j=1}^n S^j(\lambda) c_j .
\end{equation}
% explain here: reference spectrum/differential (s)cd's/doas approximations

The key idea of the DOAS method is to separate broad and narrow band
spectral structures of the absorption spectra in order to isolate the
narrow trace gas absorption features. In order to do this, some
approximations are made:
\begin{enumerate}
\item In the case where the photon path is not defined (scattered light
  measurements), the mean path followed by the photons through the
 atmosphere up to the instrument is considered;
\item The absorption cross sections are supposed to be independent of
  temperature and pressure, which allows us to introduce the concept
  of \glspl{scd};
\item Broadband variations, such as loss and gain from scattering and
  reflections by clouds and/or at the earth surface, are approximated
  by a common low order polynomial.
\end{enumerate}
Molecular absorption cross sections are fitted to the logarithm of the
ratio of the measured spectrum and the reference spectrum (i.e.\ an
extraterrestrial irradiance spectrum for satellite measurements, or a
spectrum measured around the local noon when the light path is minimum
for ground-based measurements). The resulting fit coefficients are the
integrated number of molecules per unit area along the atmospheric
light path for each trace gas, the differential \gls{scd}. The slant column depends
on the observation geometry, the position of the sun and also on
parameters such as the presence of clouds, aerosol load and surface
reflectance.

\section{DOAS retrieval}\label{sec:doas-retrieval}
\subsection{DOAS fitting}\label{sec:doas-fitting}
In the DOAS analysis, high frequency spectral structures of the
various absorbing species are used to resolve the corresponding
contributions to the measured optical density. This is achieved using
a least squares fitting procedure for the differential slant column
densities of the various species.  Large band contributions to the
atmospheric attenuation (Rayleigh and Mie scattering) are accounted
for by including a low order polynomial $P(\lambda) = \sum b_k
\lambda^k$ in the fit. Other effects, such as the Ring effect or
instrument undersampling can be treated as pseudo-absorbers.

According to eq.\ \eqref{eq:optical_density}, the DOAS retrieval is a
linear problem. This linearity is unfortunately broken down by the
need to account for additional effects, namely:
\begin{enumerate}
\item small wavelength shifts $\Delta(\lambda)$ between $I$ and $I_0$
  spectra must be corrected using appropriate shift and stretch
  parameters;
\item possible instrumental and/or atmospheric stray light or residual
  dark current signal require the introduction of an offset parameter.
\end{enumerate}
Including these necessary non-linear corrections and the polynomial
component $P(\lambda)$, we obtain the modified equation
\begin{equation}\label{eq:doas_non_linear}
\ln\left(
  \frac{I(\lambda-\Delta(\lambda))-\mathrm{offset}(\lambda)}{I_0(\lambda)}
  \right)+\sum_{j=1}^nS^j(\lambda)c_j +P(\lambda) = 0\,.
\end{equation}

\subsection{Marquardt-Levenberg Algorithm}\label{sec:marq-levenb-algor}
The \glspl{scd} $c_j$ can be retrieved by performing a least squares
fit of the measurement data to equation (\ref{eq:doas_non_linear}).
Due to the shift, stretch and offset parameters, this is a non-linear
least squares problem, which can typically be solved using a \gls{ml}
algorithm (see~\citet{Marquardt1963,Bevington1969}).  Starting from
given initial values, this algorithm searches the parameter space
using a combination of Gauss-Newton and gradient descent steps.  The
\gls{ml} algorithm is an efficient minimization algorithm for a wide
range of non-linear functions, but it may converge to a local minimum
if the initial values are too far from the global minimum.  In QDOAS,
this problem can occur during the wavelength calibration, when the
initial calibration is not accurate enough (see
\nameref{sec:calibration-problems} in section
\ref{sec:analysis-problems} of the appendix).

% explain offset modelling?
In QDOAS, the M-L algorithm is used to find the set of parameters that
minimizes the weighted\footnote{Weighting the residuals by the
  instrumental errors $\sigma^i$ is optional in QDOAS.  When no
  measurement uncertainties are used (or no error estimates are
  available), all uncertainties in equation~\eqref{eq:sum_of_squares}
  are set to $\sigma^i=1$, giving all measurement points equal weight
  in the fit.  See analysis properties in section
  \ref{sec:projects_analysis_page}} sum of squares
%\begin{equation}\label{eq:sum_of_squares}
%  S = \sum_{i=1}^M \left(\frac{Y^i - Y_{\mathrm{fit}}^i}{\sigma^i}\right)^2,
%\end{equation}
\begin{align}
  F(\vec{\alpha}) = & \frac{1}{2} \sum_{i=1}^M\left(
    \frac{f^i(\vec{\alpha})}{\sigma^i}
    \right)^2 \nonumber \\
    = &\frac{1}{2} \sum_{i=1}^M
  \left(
    \frac{\ln(I^i(\vec{a}))+\sum_{j=1}^{n}S^{ij} c_j + \sum_{k=0}^d
      b_k (\lambda_i)^k -\ln(I^i_0)}{\sigma^i}
  \right)^2,  \label{eq:sum_of_squares}
\end{align}
where
\begin{description}
\item[$\vec{\alpha}=(\vec{a},\vec{b},\vec{c})$] is the vector
  containing all fitted parameters;
\item[$\vec{a}$] is the set of parameters describing the shift,
  stretch and offset of the measured spectrum;
\item[$b_k,\, k=1,\ldots,d$] are the fitted polynomial coefficients;
\item[$c_j,\, j=1,\ldots,n$] are the fitted \glspl{scd};
\item[$\lambda_i,\, i=1,\ldots,M$] is the wavelength grid of the
  reference spectrum $I_0$, calibrated with respect to a high
  resolution solar spectrum (see section \ref{sec:wavel-calibr}).
\item[$I^i(\vec{a})$] is the measured spectrum, including shift,
  stretch and offset corrections, interpolated at wavelength
  $\lambda_i$;
\item[$S^{ij}$] is the absorption cross section of absorber species
  $j$, as measured in a laboratory, interpolated at wavelength
  $\lambda_i$;
\item[$I^i_0$] is the reference spectrum at wavelength $\lambda_i$ and
\item[$\sigma^i$] is the standard error on the measurement at
  wavelength $\lambda_i$.
\end{description}
In each iteration, the fitting algorithm will choose new values for
the parameters using the Marquardt-Levenberg algorithm, and recompute
the sum of squares $F$.  The algorithm assumes that it has converged
to a solution when the difference between two succeeding values of $F$
is smaller than a fixed small number $\epsilon$.  The convergence
criterion $\epsilon$ can be chosen by the user (see
\nameref{sec:conv-crit} in section \ref{sec:projects-properties}).

For a given value of the non-linear parameters $\vec{a}$, the
determination of the parameters $c_j$ and the polynomial coefficients
$b_k$ is a linear least squares problem.  Explicitly, we can rewrite
equation~\eqref{eq:sum_of_squares} as
\begin{equation}
  \label{eq:doas_linear}
  F = \frac{1}{2}\sum_{i=1}^M\left(
    \frac{\ln(I^i(\vec{a})/I^i_0) +  (A \cdot \vec{x})^i}{\sigma^i}
    \right)^2\,,
\end{equation}
where the matrix $A$ contains the polynomial basis $(\lambda_i)^k,
k=0,\ldots,d$ and the absorption cross sections $S^{ij}, j=1,\ldots,n$
and $\vec{x}$ represents the combined linear parameters $\vec{b}$ and
$\vec{c}$:
\begin{align}
  \label{eq:doas_linear_matrix}
  A^{il}  & = \left( (\lambda_i)^d \quad S^{ij} \right)\,,\\
  \vec{x} & = \left( \begin{array}{c}
      \vec{b}\\
      \vec{c}
    \end{array}\right)\,.
\end{align}
% describe linear equation here?
This fact is exploited in QDOAS to limit the parameter space of the
\gls{ml} algorithm to the non-linear parameters: once a new set of
values $\vec{a}'$ for the non-linear parameters is chosen, the linear
parameters $b_k$ and $c_j$ are updated using a linear least squares
algorithm, minimizing the sum of squares \eqref{eq:sum_of_squares} for
the given values of the non-linear parameters.

In QDOAS, the linear least squares problem is solved using the QR
decomposition of the matrix $A$ defined in equation
\eqref{eq:doas_linear_matrix}.

QDOAS can include further parameters in the fits
(\ref{eq:sum_of_squares}), which we omitted from the previous
description for the purpose of clarity.  Such parameters include the
width of the instrument's slit function and wavelength shifts in the
reference absorption cross sections $S^j(\lambda)$.  Mathematically,
they play the same role as the shift, stretch and offset parameters
$\vec{a}$.  Section \ref{sec:analysis_config_param} describes the
available fitting options in QDOAS and their
\index{configuration}configuration.

\subsection{Intensity fitting}\label{sec:intensity-fitting}
QDOAS also supports the so-called intensity fitting (or direct
fitting) method where measured intensities are directly fitted instead
of their logarithms. The equation used in the least squares fitting
procedure is then
\begin{equation}
  \label{eq:intensity_fitting}
  \frac{I(\lambda -\Delta(\lambda)) - \mathrm{offset}(\lambda)}{ I_0(\lambda)} - P \exp(-\sum_{j=1}^n S^jc_j) = 0\,,
\end{equation}
This method involves a decomposition of equation
(\ref{eq:intensity_fitting}) in its linear and non-linear parts:
column densities are fitted non-linearly, but the polynomial, which is
taken out of the exponential function using a Taylor expansion, and
offset are linear parameters.

Intensity fitting is sometimes preferred to optical density fitting,
for example when a poor signal leads to numerical problems when taking
the logarithms of the intensity ratio.

\subsection{Errors On Slant Column Densities}\label{sec:errors-slant-column}
% TODO: remark that errors sigma are errors on log(I) in case of
% optical density fitting, and errors on I in case of intensity
% fitting?
Uncertainties on the retrieved slant columns $c_j$ depend on:
 \begin{itemize}
 \item the sensitivity of the sum of squares $F$
   (\ref{eq:sum_of_squares}) with respect to variations of the fitted
   parameters around the minimum and
 \item the noise on the measurements.
 \end{itemize}

 Formally, the covariance matrix of the fitted parameters of a least
 squares estimate may be estimated by the inverse of the Hessian of
 the sum of squares $F$, evaluated at the fitted values of the
 parameters
\begin{equation}
  \label{eq:slnterr_general}
  \Sigma_{\vec{\alpha}} = H^{-1}\,,
\end{equation}
where $H$ is given by
\begin{equation}
  \label{eq:hessian}
  H_{kl} = \frac{\partial^2 F}{\partial \alpha_k \partial \alpha_l}\,.
\end{equation}

The actual approach used to calculate $H$ in QDOAS is different in
optical density fitting mode and in intensity fitting mode.  In
optical density fitting mode, the slant column density is fitted
linearly according to equation~\eqref{eq:doas_linear} and we use the
following expression for the covariance matrix $\Sigma_{\vec{x}}$ of
the weighted linear least squares parameter estimate:
\begin{equation}
  \label{eq:slnterr_linear}
  \Sigma_{\vec{x}} = \left(A^T W^2 A\right)^{-1}\,,
\end{equation}
where the matrix $A$ contains the linear components of the fit, as
described in equation~\eqref{eq:doas_linear_matrix} and the $(M \times
M)$ diagonal matrix $W$ contains the measurement errors
\begin{equation}
  \label{eq:weight_matrix}
  W_{ij} = \frac{1}{\sigma^i}\quad \mathrm{if}\, i=j\,.
\end{equation}
Equation \eqref{eq:slnterr_linear} implies that uncertainties on
estimated values of the non-linear parameters are not taken into
account in the reported errors on the slant columns.

When the individual measurement errors $\sigma^i$ are not available
(or the user has chosen not to weight the fit by the instrumental
errors), the weight matrix $W$ is just the identity matrix.  In this
case, the mean squared error on the measurements $\sigma^2$ may be
estimated by the reduced $\chi^2$, e.g.\ the sum of squares of the
residuals divided by the number of degrees of freedom in the fit
\begin{equation}
  \label{eq:chi_square}
  \chi^2 = \frac{\sum_{i=1}^M (f^i)^2 }{M-N}\,,
\end{equation}
where $M$ is the number of wavelengths included in the fit, and $N$ is
the total number of fitted parameters.
Equation~\eqref{eq:slnterr_linear} then becomes
\begin{equation}
  \label{eq:slnterr_linear_noweighting}
  \Sigma_{\vec{x}} = \chi^2\left(A^T A\right)^{-1}\,.
\end{equation}

In intensity fitting mode, the slant column densities are determined
using non-linear least squares fitting, and the Hessian is
approximated by the square of the Jacobian of the fitting functions
$f^i$
\begin{equation}
  \label{eq:hessian_approx}
  H_{kl} \simeq \sum_{j=1}^M \frac{\partial f^j}{\partial \alpha_k}\frac{\partial f^j}{\partial \alpha_l}\frac{1}{(\sigma^j)^2}\,.
\end{equation}
This implies that, in intensity fitting mode, uncertainties on the
linear parameters are not taken in to account in the reported slant
column errors.  As in the linear case, when the instrumental errors
are not used or not available, the $\sigma^j$ in equation
\eqref{eq:hessian_approx} are replaced by 1, and the measurement error
is approximated using the reduced $\chi^2$ of the fit instead.
Equation~\eqref{eq:slnterr_general} then becomes
\begin{equation}
  \label{eq:slnterr_nonlinear_noweighting}
  \Sigma_{\vec{\alpha}} = \chi^2 H^{-1}\,.
\end{equation}
% in practice, there are even some L-M modifications to this formula.

\subsubsection{Limitations}
The estimates for the uncertainties of the fitted parameters rely on
the statistical model used for the DOAS retrieval.  The model assumes
that the errors on the measurements at each frequency are independent
and normally distributed.  When measurement errors are correlated
or contain systematic components, or in case of fitting errors such
as a wrong fit for shift and stretch parameters, or relevant
absorbers that are not included in the fit, the uncertainties
calculated by QDOAS will underestimate the true error on the slant
column density.  Remaining residual structures after the fit, or a
value of $\chi^2 \gg 1$ when measurement errors $\sigma^i$ are taken
into account, are an indication of such a bad fit.  A discussion can
be found in~\citet{stutz_1996}.

We also note that the errors in QDOAS do not take into account the
uncertainties on the cross sections. A way to achieve this is given
in~\citet{Theys2007}.

\section{Wavelength calibration}\label{sec:wavel-calibr}
The quality of the retrieval fit strongly depends on a perfect
alignment between the spectrum to analyze, the reference spectrum and
the cross sections. The wavelength-pixels relation
(see~\citet{platt_2008}) of the % or just give expression here?
reference spectrum, previously determined in laboratory using a lamp
spectrum % what's the laboratory lamp spectrum doing here?
can be corrected using a procedure based on the alignment of the
Fraunhofer structures of the reference spectrum $I_0$ with those of an
accurately calibrated high-resolution solar reference atlas, degraded
at the resolution of the instrument, i.e. convolved with the
instrumental slit function. The reference atlas used for this purpose
is usually the~\citet{Chance2010} spectrum.

During the calibration process, the instrumental slit function can
also be characterized by repeatedly convolving the highly resolved
solar atlas with the slit function and adjusting the parameters until
the best match with the reference spectrum is found.  QDOAS allows
fitting the parameters of different line shapes, as well as their
wavelength dependence.  This is useful when the slit function provided
with the instrument is not described precisely enough in the
wavelength interval used for the retrieval.  In the same way, to
account for the moderate resolution of satellite or ground-based
instruments (about a few tenths of a nanometer), the absorption cross
sections of the trace gases have to be convolved with the instrumental
slit function (and interpolated on the final $I_0$ wavelength grid). A
good knowledge of the instrumental slit function and its potential
wavelength variation is important to avoid systematic errors in the
retrieved slant columns due to spectral shape mismatch between the
reference and atmospheric spectra~\citep{vanroozendael2002}.

\subsection{Fit of shift and stretch}\label{sec:fit-shift-stretch}
Shift and stretch (see \nameref{sec:shift} on page
\pageref{sec:shift}) can be taken into account in the wavelength
calibration scheme. To this end, the spectral interval is divided into
a number of sub-intervals.  By default, QDOAS will split the complete
calibration interval into equally sized disjoint sub-intervals, but
more advanced options are available, too (see
\nameref{sec:calibration-windows} on page
\pageref{sec:calibration-windows}).  The fitting algorithm used for
the DOAS retrieval is then applied in each sub-interval to fit the
measured intensities to those of the high-resolution solar spectrum,
according to the equation
\begin{equation}
  \label{eq:calibration_shift}
  I_0(\lambda)= I_S(\lambda - \Delta_i - a_i(\lambda-\lambda_{\mathrm{c},i}))\exp(-\sum_{j=1}^mS^jc_j)\,,
\end{equation}
where $I_S$ is the solar spectrum convolved at the resolution of the
instrument, $\Delta_i$ is a fitted constant shift in sub-interval $i$,
$\lambda_{\mathrm{c},i}$ the central wavelength of sub-interval $i$,
$a_i$ the interval's fitted stretch factor, and the $c_j$ are optional
absorber coefficients accounting for possible light absorption in the
reference spectrum $I_0$.  Typically, the calibration uses more then
one sub-interval, and we fit a polynomial through the pairs
$(\lambda_{\mathrm{c},i}, \Delta_i)$ in order to reconstruct an
accurate wavelength calibration $\Delta(\lambda)$ for the complete
analysis interval.  When the calibration is set to use only a single
interval, the shift and stretch values from this interval are used to
calculate the optimized wavelength grid for the analysis.

Several tests may be needed to determine the best configuration. One
has to find the right compromise between using enough sub-intervals to
represent the wavelength variation of the shift, and having enough
spectral information in each sub-window.  The interval calibration
should cover all analysis spectral windows.

The algorithm can take into account molecular absorption and offset
correction.  In order to help the algorithm to converge to the correct
solution, it is important to start with a preliminary wavelength grid
close enough to the real one and to limit the number of fitting
parameters (for example by fixing the value of some parameters,
typically the concentration of the fitted molecules after a first
estimation).  Another option to help the algorithm converge is to
determine the initial shift by testing for correlation:

\subsection{Calculation of the preshift}
When the wavelength shift between the high-resolution solar atlas and
the initial calibration of the reference spectrum is too large, the
fitting algorithm used in the calibration procedure may not converge
to the proper value.  If you find that this is the case, you can
choose to perform an extra ``preshift'' calculation.  With this
option, QDOAS performs the following extra steps before running the
wavelength calibration algorithm:

\begin{enumerate}
\item The high-resolution solar spectrum is degraded to instrument
  resolution using the chosen slit function settings.  If slit
  function parameters are also fit during the calibration, their
  initial values are used for this step.
\item The convolved solar spectrum and the reference spectrum are
  smoothed and normalized, with an averaging bandwidth of about 5nm.
\item Using the Nelder-Mead optimization algorithm, QDOAS determines
  the shift value resulting in the highest correlation between the
  reference and high-resolution solar spectrum.
\end{enumerate}

After this preshift step, QDOAS performs the regular wavelength
calibration on the original spectra, using the preshift value as the
initial value for the calibration shift in each sub-interval.

\subsection{Ref2/Ref1 Alignment}\label{sec:ref2ref1-alignment}
If two reference spectra are given (see section
\ref{sec:projects-properties} and figure
\ref{projects_properties_calibration}), the wavelength calibration
described in the previous section is applied to the first one.  The
shift between both reference spectra is then determined using a least
squares fit.  This shift is then applied to the reference absorption
cross sections in order to align them on the second reference
spectrum.  The measured spectra are then analyzed with respect to the
second reference spectrum.

% in detail:
% pixels are chosen based on pixel/wavelength relationship from corrected reference 1 grid,
% same pixels are selected from reference 2
% shift between ref1/ref2 is fitted
% this shift is applied to the cross sections
%
% motivation for all this?
% possibility to take solar spectrum by instrument & earthshine
% spectrum, and calibrate in two steps?
% why don't we correct the wavelength grid of reference 2?

\subsection{Characterization Of The Line
  Shape}\label{sec:char-line-shape}
If the slit function varies strongly in the fitting window and is
symmetric enough to be approached by an analytical function like a
Gaussian, it can be fitted together with the shift during the
calibration procedure.  A high resolution solar spectrum is needed.
The wavelength variation of the line shape can be accounted for in
order to convolve high-resolution cross sections before the
analysis. % should describe this more clearly... we mean varying width
          % across spectral interval?
The equation \eqref{eq:calibration_shift} used in the calibration is
then modified as follows:
\begin{equation}
  \label{eq:calibration_shift_lineshape}
  I_0(\lambda) = (L_{a_i}\ast I_S)(\lambda-\Delta_i)
  \exp\left(-\sum_{j=1}^m c_j S^j\right)\,,
\end{equation}
where $L_{a_i}(\lambda)$ is the line shape with \gls{fwhm}
parametrized by $a_i$.  QDOAS can fit the parameters for a number of
different analytical line shapes, which we describe in more detail in
the \nameref{sec:analyt-line-shap} section on page
\pageref{sec:analyt-line-shap}.

\section{Fit parameters}\label{sec:fit-parameters}
\subsection{Absorption cross sections}\label{sec:absorpt-cross-sect}
In principle, DOAS measurements are applicable to all gases with
suitable narrow absorption bands in the UV, visible or near IR
regions. However, the generally low concentrations of these compounds
in the atmosphere and the limited signal-to-noise ratio of the
spectrometers restrict the number of trace gases that can be detected.
Figure \ref{fig:doas_cross_sections} shows the absorption cross
sections of a number of trace gases that are regularly measured using
the DOAS technique.

\begin{figure}
  \centering
  \includegraphics[width=0.9\textwidth]{Algo_Cross.jpg}
  \caption{Absorption cross sections of various absorbers retrieved with
the DOAS method (figure taken from~\citet{platt_2008}).}
\label{fig:doas_cross_sections}
\end{figure}

Many spectral regions contain a large number of interfering
absorbers. In principle, because of their unique spectral structure, a
separation of the absorption is possible.  However, correlations
between absorber cross sections sometimes lead to systematic biases in
the retrieved slant columns.  In general, the correlation between
cross sections decreases if the wavelength interval is extended, but
doing so is at odds with the assumption that a single effective light
path for the entire wavelength interval is defined, leading to
systematic misfit effects that may also introduce biases in the
retrieved slant columns.  To optimize DOAS retrieval settings, a
trade-off has to be found between these effects.  A basic limitation
of the classical DOAS technique is the assumption that the atmosphere
is optically thin in the wavelength region of interest.  At shorter
wavelengths, the usable spectral range of DOAS is limited by rapidly
increasing Rayleigh scattering and \ozone{} absorption.  In addition,
‘line-absorbers’ such as \water, \oxygen, CO, \CO2 and \CH4 usually
cannot be retrieved precisely by standard DOAS algorithms because
their strong absorption also depends on pressure and temperature.

The selection of the spectral analysis window determines which
absorbers have to be included in the fitting procedure. The pressure
dependence of absorption cross sections can be neglected in UV-Visible
region. However, the temperature dependence of cross sections can be
significant (for example that of \ozone{} and \NO2). This can be
corrected in a first approximation by introducing correction factors
during the \gls{amf} calculation~\citep{Boersma2004} or,
assuming a linear dependence on temperature, by fitting two absorption
cross sections at different temperatures as described
in~\citet{vanroozendael2002}.

\subsubsection{Shift}\label{sec:shift}
Before analysis, absorption cross sections are interpolated on the
final grid of the reference spectrum which may be determined by the
program itself (automatic reference selection mode). Moreover, shift
and stretch parameters can be fitted during the processing of a
spectrum in order to obtain the best match of the absorption
structures.

\subsubsection{Interpolation/Convolution}
Cross sections can be pre-convolved and interpolated\footnote{By
  default, QDOAS uses cubic spline interpolation.} on an appropriate
wavelength grid prior to the analysis. However, the direct use of
high-resolution cross sections which can be convolved in real-time
with a predefined slit function or with the slit function determined
by the wavelength calibration procedure, is more convenient.

\subsubsection{Differential cross sections}
The aim of calculating differential absorption cross sections is to
separate narrow spectral features from unstructured absorption not
useful in the DOAS method.  Figure
\ref{fig:differential_cross_section} gives an illustration.

\begin{figure}
  \centering
  \begin{tabular}{cc}
    \includegraphics[width=0.45\textwidth]{Algo_Cross_1.jpg}
    & \includegraphics[width=0.45\textwidth]{Algo_Cross_2.jpg}
  \end{tabular}
  \caption{Cross section $\mathrm{O}_3$ and optical thickness $\tau$ are separated into a narrow ($\mathrm{O}_3'$ and $\tau'$)
and a broad band part ($\mathrm{O}_{3p}$ and $\tau_p$) by
an adequate filtering procedure.}
\label{fig:differential_cross_section}
\end{figure}

Differential cross sections can be calculated by orthonormalization
with respect to an orthogonal set formed with the low order component
vectors (usually a base of order 2 or 3) of the polynomial.  The
vectors are orthogonalized using the Gram-Schmidt
algorithm.  This method gives also the possibility to separate the
spectral structures of temperature dependent cross sections by
mutually orthogonalizing them.  Differential cross sections can
alternatively be obtained by high-pass filtering: a filter is applied
iteratively on the cross section and the resulting cross section is
subtracted from the original one. High-pass filtering is presently
only supported in optical density fitting mode (DOAS mode).  Note that
it is not necessary to calculate differential cross sections when the
baseline is around 0 such as for \water, \tetraox{} or \H2CO.

\subsubsection{Low-pass filtering}
Low-pass filters can be applied to both spectra and absorption cross
sections. A large choice of filters is proposed (see `Filtering
Page', section \ref{sec:projects_filtering_page}).

\subsubsection{Wavelength-dependent AMF}
Absorption cross sections can be replaced by geometrically corrected
cross sections that take into account the wavelength dependency of the
\gls{amf}. The correction is based on the equation
\begin{equation}
  \label{eq:wavelength_dep_amf}
  \ln(I_0)-\ln(I) = S(\lambda) \mathrm{AMF}(\lambda,\mathrm{SZA})\,,
\end{equation}
where $\mathrm{AMF}(\lambda,\mathrm{SZA})$ is the AMF calculated for a
given wavelength $\lambda$ and a given \gls{sza}.
% fitted density is then already the vertical column density?

\subsection{Polynomial}\label{sec:polynomial}
In the DOAS technique, absorption cross sections of the considered
mole\-cules are highly structured, while scattering by molecules and
particles (Rayleigh and Mie scattering), as well as reflection at the
surface, have broadband dependencies that can be approximated by a low
order polynomial. In QDOAS, polynomials up to degree 5 can be
fitted. The low order components of the polynomial can be used to
calculate differential cross sections when this is preferred to
high-pass filtering.

\subsection{Shift and Stretch}\label{sec:shift-stretch}
Shift and stretch parameters allow to correct possible misalignment
between the various wavelength-dependent quantities involved in the
data evaluation (i.e.\ measured and reference spectra as well as
absorption cross sections).  Shift and stretch parameters may be
fitted or simply applied to any wavelength-dependent quantity,
according to the equation
\begin{equation}
  \label{eq:shift_stretch}
  \Delta\lambda = a + b(\lambda-\lambda_0) + c(\lambda-\lambda_0)^2\, ,
\end{equation}
where $\lambda$ is the wavelength according to the original
calibration and $\lambda_0$ is the center wavelength of the current
spectral range.  The parameters $a$,$b$ and $c$ describe the offset
and the first and second order stretch applied to $\lambda$.

\subsection{Offset Correction}\label{sec:offset-correction}
An ideal spectrometer in an ideal atmosphere would measure the part of
the sunlight that has been elastically scattered by air molecules and
particles.  In a real experiment, however, a number of possible
additional sources may contribute to the measured intensity, adding an
“offset” to the ideal Rayleigh/Mie contribution.  In addition to the
Ring effect, which is, to a first approximation, a natural source of
offset, one must also account for instrumental sources of offset, such
as stray light in the spectrometer and dark current of the detector.
The offset component in equation~\eqref{eq:doas_non_linear} accounts
for these effects.  QDOAS models the offset using a polynomial
\begin{equation}
  \label{eq:offset}
  \mathrm{offset}(\lambda) = (a + b(\lambda-\lambda_0) + c(\lambda-\lambda_0)^2) \bar{I}\, ,
\end{equation}
where $\lambda_0$ is the center wavelength of the spectral analysis
window, $\bar{I}$ is the average intensity and $a$, $b$ and $c$ are
fitted parameters.

Due to the normalization by $\bar{I}$, the offset values can be easily
interpreted relatively to the absolute intensity of the spectrum
(percent offset).  In DOAS, the offset is usually fitted as a
non-linear parameter.  Sometimes, this could lead to numerical errors
in the evaluation of the logarithm, typically in the near UV region
where the level of the signal is very low.  In this case, we can fit
the offset as a linear parameter.  Expanding the left-hand side of
equation~\eqref{eq:doas_non_linear} to first order, we obtain
\begin{equation}
  \label{eq:offset_linear}
  \ln(I(\lambda -\Delta(\lambda))) = \ln(I_0(\lambda))-\tau(\lambda) + \frac{\mathrm{offset}(\lambda)}{I(\lambda+\Delta\lambda)}\,.
\end{equation}
 % why not also normalize with respect to average intensity in this case?

\subsection{Ring Effect}\label{sec:ring-effect}
Because of \gls{rrs}, a small fraction of the
incident photons undergo a wavelength change of a few nanometres,
i.e.\ a part of the scattering is inelastic. This causes an intensity
loss at their incident wavelength and a gain at the neighbouring
wavelengths to which they are redistributed.  RRS causes the so-called
“filling-in” of Fraunhofer lines, which have a slightly different
shape in the earthshine radiance than in the direct solar light.  This
effect was first discovered by~\citet{grainger_1962} and is referred
to as the Ring effect.  The atmospheric absorption lines are also
broadened by RRS events occurring after absorption (molecular Ring
effect).  Although RRS accounts for only a few percent of the measured
intensity, it significantly affects DOAS measurements of scattered
radiation since typical trace gas absorptions are of the order of a
percent or less.  If not properly corrected, the Ring effect produces
strongly structured residuals in the differential optical density, due
to the fact that Fraunhofer lines do not cancel perfectly between $I$
and $I_0$. Especially in the UV spectral range, the remaining spectral
structures can by far exceed the structures of weak atmospheric
absorbers.

Usually the Ring effect is taken into account by including an
additional ‘absorber’ in the DOAS fit. Ring cross sections
$S_{\mathrm{Ring}}(\lambda)$ can be
measured~\citep{Solomon_Schmeltekopf_Sanders_1987} or
calculated~\citep{Chance1997}. The Ring effect can be approximated
using the following development for an optically thin
atmosphere~\citep{Wagner1999,vanroozendael2002}:

One can consider that in any scattered light observation, the light
detected by the satellite instrument ($I_\mathrm{meas}$) is the sum of
elastic and inelastic scattering processes:
\begin{equation}
  \label{eq:rrs_i_meas}
  I_\mathrm{meas} = I_{\mathrm{elastic}} + I_{\mathrm{RRS}} = I_0 e^{-\tau} + I_{\mathrm{RRS}}\,.
\end{equation}
To analyze a measured spectrum, the logarithm of $I_{\mathrm{meas}}$
is taken. Since $I_{\mathrm{RRS}}$ is very small compared to
$I_{\mathrm{elastic}}$, the logarithm can be approximated by the first
two terms of the Taylor expansion:
\begin{equation}
  \label{eq:rrs_expand_i_meas}
  \ln(I_\mathrm{meas}) = \ln(I_{\mathrm{elastic}}) + \frac{I_{\mathrm{RRS}}}{I_{\mathrm{elastic}}}
  = \ln I_0 - \tau + \frac{I_{\mathrm{RRS}}}{I_{\mathrm{elastic}}} \,.
\end{equation}
The last term $I_{\mathrm{RRS}}/I_{\mathrm{elastic}}$ is the Ring term
which can be approximated by the product of a Ring coefficient
$\alpha_{\mathrm{Ring}}$ and a Ring cross section
$S_{\mathrm{Ring}}$. $\alpha_{\mathrm{Ring}}$ can be fitted together
with other absorbers in the DOAS procedure.

A good approximation of the Ring cross section can be obtained from
calculation of a source term $I_0^{\mathrm{Raman}}$ for Raman
scattering derived by simple convolution of the solar spectrum with
Raman cross sections, i.e with calculated \N2 and \oxygen{} RRS cross
sections. The term $I_{\mathrm{RRS}}/I_{\mathrm{elastic}}$ in
equation~\eqref{eq:rrs_expand_i_meas} is then substituted by
$\alpha_{\mathrm{Ring}} S_{\mathrm{Ring}}$ (where $S_{\mathrm{Ring}}
=I_0^{\mathrm{Raman}}/I_0$), assuming that the molecular Ring effect
can be neglected~\citep{Chance1997}. The QDOAS Ring tool calculates a
Ring cross section according to this
approach~\citep{Wagner1999}. The normalization of the Raman spectrum is optional.

\subsection{Undersampling Correction}\label{sec:unders-corr}
The undersampling is a well-known problem of
GOME~\citep{Chance1998}. It arises from the poor sampling ratio of the
GOME instrument, which results in a loss of spectral information when
interpolating earthshine spectra during the DOAS fitting
process. Undersampling is only a major problem when the sampling ratio
is close to 1.  In the case of GOME, this ratio is estimated at
0.16/0.112=1.43 at 340 nm.

The undersampling artifacts can be corrected using an ad-hoc cross
section obtained by simulating the effect based on a high-resolution
solar reference~\citep{Chance1998}.  This cross section is fitted as a
linear parameter, in both intensity fitting and in optical density
fitting modes.  The procedure to calculate the correction differs
slightly between the two fitting modes:

In intensity fitting mode, oversampled
and undersampled spectra are calculated as follows:
\begin{equation}
  \label{eq:undersampling}
  U(\lambda+\Delta) = \mathrm{over}(\lambda+\Delta) - \mathrm{under}^{(\Delta)}(\lambda)\,,
\end{equation}
where $\mathrm{over}(\lambda + \Delta)$ is a high-resolution solar
spectrum convolved on its original grid and interpolated on the final
grid $\lambda + \Delta$ and $\mathrm{under}^{(\Delta)}(\lambda)$ is a
high-resolution solar spectrum convolved on grid $\lambda$ and
subsequently interpolated on the final grid $\lambda + \Delta$.
% does it make a difference if we convolve on one grid or another?

Residuals are improved by adding a “second phase” of undersampling
\begin{equation}
  \label{eq:undersampling2}
  U_2(\lambda) = \mathrm{over}(\lambda) - \mathrm{under}^{(\Delta)}(\lambda)\,.
\end{equation}
In optical density fitting, we fit the logarithm of the ratio of the
measured intensities.  We therefore use the corresponding equations
\begin{align}
  \label{eq:undersampling_doas}
  U(\lambda+\Delta) & = \ln\left(\frac{\mathrm{over}(\lambda+\Delta)}{\mathrm{under}^{(\Delta)}(\lambda)}\right), \\
  U_2(\lambda) & = \ln\left(\frac{\mathrm{over}(\lambda)}{\mathrm{under}^{(\Delta)}(\lambda)}\right).
\end{align}
Undersampling cross sections can be pre-calculated using the QDOAS
undersampling tool or they can be calculated in real time, just after
the wavelength calibration procedure using the corrected wavelength
grid and the characterized slit function.

\subsection{AMF}\label{sec:amf}
To obtain the \gls{vcd} from the fitted \gls{scd}, we have to divide
the SCD by the \gls{amf}
\begin{equation}
  \label{eq:air_mass_factor}
  \mathrm{VCD} = \frac{\mathrm{SCD}}{\mathrm{AMF}}\,.
\end{equation}
AMF's are not calculated by QDOAS but may be imported from ASCII files
(for details, refer to appendix \ref{cha:input-files-format}). QDOAS
supports solar zenith angle-dependent AMF, wavelength-dependent AMF
and climatological AMF.

\subsection{Graphical Display Of Spectral Fit Results}\label{sec:graph-displ-spectr}
According to the selected fitting method, different formulas are used
to display the measured data and the calculated fit results.  Tables
\ref{tab:display_results_od} and \ref{tab:display_results_intensity}
explain which quantities are displayed on the fit results page for
optical density fitting mode and intensity fitting mode.  The
following quantities are used:
 \begin{align}
   R_{\mathrm{od}} & = \ln\left(\frac{I_0 \cdot e^{-\sum S^j c_j - P}}{I-\mathrm{offset}}\right)\,,\label{eq:residual}\\
   R_{\mathrm{norm}} & = \frac{P\cdot I_0\cdot e^{-\sum S^j c_j}}{I-\mathrm{offset}} -1 \,,\\
   R_{\mathrm{i}} & = \ln\left(\frac{P\cdot I_0\cdot e^{-\sum S^j c_j}}{I-\mathrm{offset}}\right)\,,\\
   P_{\mathrm{meas}} &= \frac{I-\mathrm{offset}}{I_0\cdot e^{-\sum S^j c_j}}\,,\\
   O & = \ln(1-\mathrm{offset}/I)\,,\label{eq:display_offset}
 \end{align}
 where $I$ is the measured spectrum, $I_0$ is the reference spectrum,
 $P$ the fitted polynomial, $S^j$ the absorption cross section for
 absorber species $j$ and $c_j$ the corresponding fitted
 concentration.

 \begin{table}
   \centering
   \begin{tabular}{lll} \toprule
     & Calculated & Measured\\
     \midrule
     Residual & & $R_\mathrm{od}$  \\
     Cross sections & $S^j\cdot c_j$ & $ S^j\cdot c_j + R_\mathrm{od}$ \\
     Polynomial & $P$ & $P + R_\mathrm{od}$\\
     Offset & $O$ & $O + R_\mathrm{od}$ \\
     \bottomrule
   \end{tabular}
   \caption{Formulas used by QDOAS to display the results of a fit in optical density fitting mode.  $S^j$ is the reference cross section for species $j$ and $c_j$ the corresponding SCD.  $P(\lambda)$ is the polynomial obtained from the fit.  See equations \eqref{eq:residual}-\eqref{eq:display_offset} for the definitions of the quantities $R_\mathrm{od}$ and~$O$.}
   \label{tab:display_results_od}
 \end{table}


 \begin{table}
   \centering
   \begin{tabular}{lll} \toprule
     & Calculated & Measured\\
     \midrule
     Normalized Residual & & $R_\mathrm{norm}$  \\
     OD Residual & & $R_\mathrm{i}$\\
     Cross sections & $S^j\cdot c_j$ & $ S^j\cdot c_j + R_\mathrm{i}$ \\
     Polynomial & $P$ & $P_\mathrm{meas}$\\
     Offset & $O$ & $O + R_\mathrm{i}$ \\
    \bottomrule
   \end{tabular}
   \caption{Formulas used by QDOAS to display the results of a fit in intensity fitting mode.  $S^j$ is the reference cross section for species $j$ and $c_j$ the corresponding SCD.  $P(\lambda)$ is the polynomial obtained from the fit.  See equations \eqref{eq:residual}-\eqref{eq:display_offset} for the definitions of the quantities $R_\mathrm{norm}, R_\mathrm{i}$ and~$O$.}
   \label{tab:display_results_intensity}
 \end{table}


\section{Block-diagram structure of the program}
\label{sec:block-diagr-struct}
The following figures give an overview of the structure of the
program. Figure \ref{fig:overal_structure} describes the general
organization of the processor.  Figures
\ref{fig:wavelength_calibration} and \ref{fig:slant_column_fitting}
show the structure of the wavelength calibration and slant column
fitting modules, while the coupled non-linear least squares fitting
algorithm is sketched in Figure \ref{fig:least_squares}.

\begin{figure}
  \centerline{\input{img/overall_structure.pdf_t}}
  \caption{Overall structure of the spectral analysis program.}
  \label{fig:overal_structure}
\end{figure}

\begin{figure}
  \centerline{\input{img/wavelength_calibration.pdf_t}}
  \caption{Structure of the wavelength calibration module.}
  \label{fig:wavelength_calibration}
\end{figure}

\begin{figure}
  \centerline{\input{img/slant_column_fitting.pdf_t}}
  \caption{Structure of the slant column fitting module.  Cross
    section processing only takes place during the pre-processing
    phase, unless cross sections must be dynamically modified during
    the fit (e.g.\ when fitting a cross section shift).  }
  \label{fig:slant_column_fitting}
\end{figure}

\begin{figure}
  \centerline{\input{img/least_squares.pdf_t}}
  \caption{Structure of the non-linear least squares fitting algorithm
    used in the slant column fitting module and in the wavelength
    calibration module.}
  \label{fig:least_squares}
\end{figure}
\FloatBarrier

\section{Convolution}\label{sec:convolution}
\subsection{Definition}\label{sec:conv_definition}
The convolution of a spectrum $I$ by an instrumental slit function $F$
is given by the integral
\begin{equation}
  \label{eq:convolution}
  (F \ast I)(\lambda) = \int I(\lambda') F(\lambda-\lambda') \mathrm{d}\lambda'\,.
\end{equation}
In QDOAS, this integral is calculated using the trapezoidal rule. The
integration interval is defined by the width of the slit function.

The three tools (Convolution, Ring and Usamp) provided in the QDOAS
package, support convolution with one of the following functions:
\begin{itemize} \itemsep -6pt
\item the Gaussian line shape,
\item the $2n$-Lorentzian line shape,
\item the Voigt profile,
\item the error function,
\item an asymmetrical Gaussian line shape,
\item the super-Gaussian line shape (\citep{Beirle2017}),
\item a boxcar (using Fourier Transform) % describe these also in the section on slit functions/line shapes?
\item Norton-Beer strong (using Fourier Transform)
\item arbitrary line shapes provided in ASCII files.
\end{itemize}

\subsection{Analytical line shapes}\label{sec:analyt-line-shap}
QDOAS can calculate the Gaussian, Lorentzian, Voigt, asymmetrical Gaussian, super-Gaussian and
error function line shapes for different values of their parameters.
When one of these analytical line shapes is used, QDOAS can fit its
parameters during the wavelength calibration procedure (see section
\ref{sec:char-line-shape}, page \pageref{sec:char-line-shape}).

In order to make the convolution algorithm faster, analytical slit
functions are pre-calculated on a suitable wavelength grid (determined
in order to have 18 pixels at
\gls{fwhm}) % what does this mean?  does it mean that a wavelength interval of FWHM contains 18 pixels?
and then interpolated on the grid of the spectrum.  For Gaussian and
error function line shapes, a \gls{fft} algorithm is used whenever possible
to speed up the calculation (e.g.\ within the wavelength calibration
procedure).

\begin{figure}
\centering
\subfigure[]{
  \includegraphics[]{gauss}
  \label{fig:gaussian}}
\subfigure[]{
  \includegraphics[]{asym02}
  \label{fig:asym02}
}
\subfigure[]{
  \includegraphics[]{asym04}
  \label{fig:asym04}
}
\subfigure[]{
  \includegraphics[]{voigt}
  \label{fig:voigt}
}
\subfigure[]{
  \includegraphics[]{lorentz1}
  \label{fig:lorentzian}
}
\subfigure[]{
  \includegraphics[]{lorentz2}
  \label{fig:lorentzian_flat}
}
\caption{QDOAS can use different line shapes to describe the
  instrumental slit function.  The plots above illustrate the
  Gaussian~(a), asymmetrical~(b, c), Voigt~(d) and Lorentzian~(e,~f)
  profile shapes.  The asymmetry factor in (b) is 0.2, that of line
  shape (c) is 0.4.  The Lorentzian (e) is the standard Lorentzian
  ($n=1$) and the Lorentzian in (f) has $n=2$.}\label{fig:lineshapes}
\end{figure}

\textbf{As a convention, QDOAS always requests the full width at half
  maximum (FWHM) of the line shape in the user interface.  This
  parameter is represented by $\sigma$ in the formulas below.}

\subsubsection{Gaussian}
The standard expression used to approximate instrumental slit
functions is the Gaussian function (see figure
\ref{fig:gaussian}). The Gaussian is the exact line shape in the
diffraction limit, i.e. in the case of an infinitely thin entrance
slit. The normalized Gaussian is given by
\begin{equation}
  \label{eq:gaussian}
  G_\sigma(x) =  \frac{1}{a\sqrt{\pi}}\exp\left(-\frac{x^2}{a^2}\right)\,,
\end{equation}
where $a$ depends on the Gaussian full width at half-maximum $\sigma$ as follows:
\begin{equation}
  \label{eq:gaussian_fwhm}
  \sigma = 2\sqrt{\ln 2}a\,.
\end{equation}

\subsubsection{Error function}
For a large entrance slit, the slit function can be approximated by
the convolution of a boxcar and a Gaussian function.  The resulting
slit function is given by
\begin{equation}
  \label{eq:error_function_slit}
  E_a(x) = \frac{1}{4\Delta}\left(
    \mathrm{erf}\left(\frac{x+\Delta}{a}\right) - \mathrm{erf}\left(\frac{x-\Delta}{a}\right) \right)\, ,
\end{equation}
where $a$ depends on the Gaussian full width at half-maximum $\sigma$ (see \ref{eq:gaussian_fwhm}),
$\Delta$ is half the boxcar width $b$ and $\mathrm{erf}(x)$ is the
error function, defined by
\begin{equation}
  \label{eq:error_function}
  \mathrm{erf}(x) = \frac{2}{\sqrt{\pi}}\int_0^x \exp(-t^2)\mathrm{d}t\,.
\end{equation}

Full width at half maximum of the Gaussian $\sigma$ (see
\ref{eq:gaussian_fwhm}) and boxcar width $b$ (2*$\Delta$) are the two
parameters requested by QDOAS for the error function.

\subsubsection{2$n$-Lorentzian}
This line shape is a generalization of the well-known Lorentzian
function ($n=1$, see figure \ref{fig:lorentzian}).  It is useful to
approximate slit functions with a shape significantly different from
the Gaussian one. Low order 2$n$-Lorentz can mimic line shapes with
large wings (typically the standard Lorentzian profile), or line
shapes having a flat top (for higher values of $n$, see figure
\ref{fig:lorentzian_flat}). The generic expression for this function
is
\begin{equation}
  \label{eq:2nlorentz}
  L_{2n}(x) = \frac{\tau^{2n}}{x^{2n}+\tau^{2n}}\,,
\end{equation}
with $\tau=\sigma/2$ ($\sigma$ is the full width at half maximum).

Parameters requested by QDOAS are the full width at half maximum
$\sigma$ and the parameter $n$, called ``Order''.

\subsubsection{Voigt profile}
The Voigt profile function (see figure \ref{fig:voigt}) is the
convolution of a Gaussian and a Lorentzian function. This function is
used in a wide range of contexts, and the optimization of its
computation has received much attention. The Voigt profile is usually
expressed as
\begin{equation}
  \label{eq:voigt_profile}
  K_(x,y) = \frac{y}{\pi} \int_{-\infty}^{\infty} \frac{\exp(-t^2)}{(x-t)^2+y^2}\mathrm{d}t\,,
\end{equation}
where $x$ is the distance from the line center in units of Gaussian
half-widths and $y$ is the ratio of the Lorentzian width to the
Gaussion width. Numerous ways have been proposed to compute the Voigt
function. A rather efficient algorithm is described
by~\citet{Kuntz_1997}. % is this the algorithm used in QDOAS?

\subsubsection{Asymmetrical Gaussian}
To handle asymmetrical line shapes, QDOAS can fit an asymmetry factor
in addition to the \gls{fwhm}.  The line shape used is a modified version of
the Gaussian line shape (\ref{eq:gaussian})
\begin{equation}
  \label{eq:asymmetrical_line}
  G'(x) = \frac{1}{a\sqrt{\pi}}\exp\left(
    -\frac{x^2}{a^2(1+\mathrm{sign}(x)b)^2}\right)\,,
\end{equation}
where, in addition to the FWHM $\sigma$, the asymmetry factor $b$ is
now fitted as well.  For a value of $b=0$, it reduces to the ordinary
Gaussian line shape.  Figure \ref{fig:asym02} illustrates the line
shape with an asymmetry factor of $b=0.2$ and \ref{fig:asym04} gives
an illustration for $b=0.4$.

\subsubsection{Super Gaussian}
This line shape is a generalization of the Gaussian with the introduction of two independent parameters, w and k, determining
the width and shape of the slit function.  This line shape has been the subject of studies by Steffen Beirle from the Max-Planck Institute in Germany and summarized in a recent paper (\citep{Beirle2017}).
The current implementation in QDOAS is operational but not optimized in terms of performance.

\subsection{User-defined line shape}\label{sec:userdef-line-shap}

An alternative to fit asymmetrical slit functions is to provide a line shape in an ASCII file and to apply on the grid of the line shape a stretch factor that can be different for each side of the line shape.
The ``Stretch on wavelength'' file in the \textbf{Convolution tool} is a ASCII file with three columns : the wavelength calibration and the two values of the stretch to apply on the grid of the line shape for each wavelength of the detector.

Lookup tables of slit functions are also accepted.  This means that a
line shape is provided for different wavelengths.  For a specific
wavelength, a line shape is first calculated by linear
2D-interpolation in the lookup table and the fitted stretch factors
are then applied to the line shape wavelength grid before convolution.

\subsection{Wavelength dependency}
QDOAS supports convolution with wavelength-dependent slit functions:
\begin{itemize}
\item Slit function files provided by the user may contain multiple
  columns describing the shape of the slit function at various
  wavelengths.  The first line of the file should list the wavelengths
  for which the slit function is provided.  QDOAS will then generate a
  lookup table from the provided slit functions, and calculate the
  slit function at intermediate wavelength values using linear
  interpolation.
  % explain this in more detail (and perhaps move this to the tools or
  % properties chapters?)
\item For analytical line shapes, the user can provide a two-column
  ASCII file containing the wavelengths and corresponding values of
  the parameters for each requested parameter (for example, the
  Gaussian width and the boxcar width for the error function).  Files
  with the correct parameter calibration can be saved from the SFP
  plots obtained after the wavelength calibration procedure.
\end{itemize}
For more details on how to compute such a convolution using QDOAS, see
section \ref{sec:conv-tool}, page \pageref{sec:slit-function-page}.

\subsection{Correction for the Solar $I_0$-effect}\label{sec:correction-solar-i_0}
The $I_0$-effect is due to the difference between cross sections
measured in the laboratory with a smooth light source and the
atmospheric absorption measured by an instrument of lower resolution
with a structured solar light source~\citep{aliwell_2002}.  As
described in equation~\eqref{eq:optical_density}, the DOAS technique
attempts to isolate the atmospheric absorptions by calculating the
ratio of a reference spectrum $I_0$ and a measured spectrum $I$.
However, because the two spectra in this ratio have been filtered by
the instrument function before the calculation of the ratio, complete
removal of the Fraunhofer structure from the solar reference spectrum
is not possible:

If the $I_0$-effect is not accounted for, the cross sections $S$ used
in the DOAS retrieval are laboratory cross sections $S_0$ convolved
with the instrumental slit function $F$:
\begin{equation}
  \label{eq:cross_section_convolved}
  S(\lambda) = (F\ast S_0)(\lambda) = \int S_0(\lambda') F(\lambda-\lambda')\mathrm{d}\lambda' \,.
\end{equation}
The intensities measured by the spectrometer can be defined as
\begin{align}
  \label{eq:incoming_intensity}
  I_0^{\mathrm{meas}}(\lambda) &= (F\ast I_0)(\lambda)\,,\\
  I^{\mathrm{meas}}(\lambda) &= (F\ast I)(\lambda) = \left( F \ast I_0 \exp(-S_0 \cdot c)\right)(\lambda)\,.
\end{align}
Taking into account the instrumental function $F$,
equation~\eqref{eq:optical_density} becomes:
\begin{align}
  \tau = \ln \frac{I_0^{\mathrm{meas}}(\lambda)}{I^{\mathrm{meas}}(\lambda)}
  &  = \ln \frac{ (F \ast I_0)(\lambda)}{ (F \ast I_0 \exp(-S_0 \cdot c))(\lambda)}\\
  & \simeq c \cdot (F * S_0)(\lambda)\,. \label{eq:convolution_approximation}
\end{align}
Thus, the fitting of convolved laboratory cross sections to the
optical depth measured by a low resolving instrument relies on an
approximation.  The error induced by this approximation depends on the
absorption strength and the relation between the width of the slit
function and the width of the spectral structures in the light source
$I_0$ and the cross section $S_0$.  For most of the atmospheric
absorbers the $I_0$-effect is negligible but in some cases, the
approximation (\ref{eq:convolution_approximation}) is not accurate
enough. This is particularly the case for strong atmospheric absorbers
such as \ozone{}.

This effect may be dealt with to a good approximation by correcting
the laboratory cross sections with the solar Fraunhofer spectrum: for
a particular absorber, a highly resolved solar spectrum $I_0(\lambda)$
and a simulated absorption spectrum $I^*_c(\lambda)$ with a known slant
column $c$ are both convolved with the instrument slit function
\begin{align}
  I_0^*(\lambda)&= (F \ast I_0) (\lambda)\, ,\\
  I^*_c (\lambda) &= (F \ast I_0 \exp( -S_0 \cdot c)) (\lambda)\,.
\end{align}
The $I_0$-corrected cross section $S*$ is then defined by
\begin{equation}
  \label{eq:I0_corrected_cross_section}
  S^*_c(\lambda) = -\frac{1}{c}\ln \left(
    \frac{ I^*_c(\lambda) }{I^*_0(\lambda)}
  \right)\,.
\end{equation}
The $I_0$-corrected cross section derived in this manner perfectly
matches the absorptions in the measured atmospheric spectrum, provided
that the SCD $c$ used in the calculation is equal to the true
atmospheric SCD. However, this is not a critical point and the $I_0$
corrected cross sections can be used for a large range of atmospheric
slant column densities~\citep{Wagner1999}.

\subsection{Deconvolution}\label{sec:deconvolution}
The QDOAS convolution tool implements a de-convolution option. In this
case, the convolution is performed using an effective slit function
obtained after Fourier transform manipulation of specified convolution
and de-convolution functions. The effective slit function is
calculated as follows:
\begin{equation}
  \label{eq:deconvolution}
  S_{\mathrm{eff}} = \mathfrak{I}^{-1} \left[
    \frac{ \mathfrak{I}(S_1) }{\mathfrak{I}(S_2)}
    \right]\,,
\end{equation}
$\mathfrak{I}$ represents the Fourier transform, and $S_1$ and $S_2$ are
respectively the convolution and de-convolution slit functions. Great
care is taken in the algorithm to avoid noise-corruption effects when
taking the ratio of the Fourier transforms.

%%% Local Variables:
%%% mode: latex
%%% TeX-master: "QDOAS_manual"
%%% TeX-PDF-mode: t
%%% End:
